% Clase 3 --- Gráfica de la función signo (§4.2).
% Constante por trozos: -1 en R^-, 0 en {0}, 1 en R^+. Los dos segmentos NO
% alcanzan (0,-1) y (0,1) (extremos abiertos); el punto aislado (0,0) sí
% pertenece a la gráfica (extremo cerrado). El docente lo describe así en
% clase sin llegar a dibujarlo él mismo; se reconstruye la descripción.
\begin{center}
\begin{tikzpicture}
\begin{axis}[
    name=signoaxis,
    calcfig,
    xmin=-2.6, xmax=2.6,
    ymin=-1.6, ymax=1.6,
    xtick={-2,-1,1,2},
    ytick={-1,1},
    xlabel={$x$},
    ylabel={$\mathrm{sgn}(x)$},
  ]
    % rama negativa: x < 0, altura -1 (extremo abierto en x=0)
    \addplot[figblue, line width=1.1pt, domain=-2.4:-0.02, samples=2] {-1};
    % rama positiva: x > 0, altura 1 (extremo abierto en x=0)
    \addplot[figblue, line width=1.1pt, domain=0.02:2.4, samples=2] {1};
    % punto aislado (0,0): pertenece a la gráfica -> círculo lleno
    \node[figptocerrado] at (axis cs:0,0) {};
    % extremos abiertos (0,-1) y (0,1): no pertenecen a la gráfica -> círculo vacío
    \node[figptoabierto] at (axis cs:0,-1) {};
    \node[figptoabierto] at (axis cs:0,1) {};
  \end{axis}
\end{tikzpicture}\\[0.5em]
{\small\itshape La función signo es constante por trozos: vale $-1$ en todo
$x<0$, vale $1$ en todo $x>0$, y el salto en $x=0$ deja un único punto
$(0,0)$ en la gráfica (círculo lleno). Los extremos $(0,-1)$ y $(0,1)$ de los
dos segmentos horizontales \textbf{no} pertenecen a la gráfica (círculos
vacíos): la función no toma esos valores en $x=0$.}
\end{center}
